\documentclass[11pt]{article}
\usepackage[margin=1in]{geometry}
\usepackage{amsmath}
\usepackage{booktabs}
\usepackage{enumitem}
\usepackage{graphicx}
\usepackage{array}
\usepackage{caption}
\usepackage{multicol}
\usepackage{tikz}
\usepackage{import}
\usepackage{adjustbox}
\setlist{nosep}

\import{../}{labflowchart.sty}

\begin{document}

\begin{center}
  {\LARGE \textbf{SCH4U DIY Thesis Lab}}\\[1em]
  {\large \textbf{A Non-Standard Vinegar Cell and the Nernst Equation}}\\[0.5em]
  {\normalsize Name: Aaron Huang \hspace{2cm} Date: 5/29/2026}
\end{center}

\section*{Purpose}
\begin{itemize}
  \item To compare the measured voltage of non-standard vinegar-based electrochemical cells using a copper cathode against the theoretical cell potentials predicted by the Nernst equation.
  \item To calculate the pH of the vinegar electrolyte from experimental cell potentials and compare these results to direct pH measurements.
\end{itemize}

\section*{Introduction}
Electrochemical cells convert chemical energy from spontaneous redox reactions into electrical energy. In this non-standard cell, hydrogen ions from an acetic acid solution (vinegar) undergo reduction at an inert copper cathode surface, while various metals (zinc, aluminum, or tin) undergo oxidation at the anode. Because the hydronium ion concentration in weak acetic acid is significantly lower than the standard $1.0\text{ M}$, the resulting cell potential is lower than the standard reduction potential, a relationship modeled quantitatively using the Nernst equation.

The experimental procedure involved preparing various metal anode strips and pairing each with a common copper cathode in a beaker filled with white vinegar. A multimeter was utilized to measure the resulting cell potentials, and a digital pH meter was used to record the baseline pH of the vinegar electrolyte. These measurements allowed for a direct mathematical comparison between the theoretical and experimental cell potentials.

Safety goggles were used to protect against splashing of vinegar into eyes.

\section*{Materials}
\begin{itemize}
  \item 100 mL beaker
  \item Copper strip (electrode)
  \item Zinc strip (electrode)
  \item Aluminum strip (electrode)
  \item Tin strip (electrode)
  \item White vinegar (acetic acid solution)
  \item Multimeter with leads
  \item Alligator clips
  \item Distilled water
  \item Safety goggles
\end{itemize}

\section*{Procedure Flowchart}
\begin{procedureflowchart}
    \lstep{Electrodes}
    \rstep{Rinse each strip with water and dry with a paper towel}
    \lstep{Vinegar}
    \rstep{Pour to around 1/4 full in the beaker}
    \lstep{pH Meter}
    \bstep
    \rstep{Rinse the meter in distilled water}
    \bstep
    \rstep{Dry the pH meter probe}
    \rstep{Measure and record the pH of the vinegar}
    \lstep{Multimeter}
    \rstep{Clip the leads to the strips}
    \rstep{Set to Voltage}
    \lstep{Electrodes}
    \lstep{Prepared Vinegar}
    \rstep{Place the copper in beaker}
    \rstep{Place the zinc in beaker}
    \rstep{Record Cell Voltage}
    \lstep{Aluminum and Tin Strips}
    \rstep{Repeat the process for aluminum and tin strips as the anode}
\end{procedureflowchart}
\newpage
\section*{Observation}
% table 1
\begin{table}[ht]
  \caption{Voltage and Qualitative Measurements}
\begin{tabular}{p{4cm} p{3.5cm} p{5cm}}
\toprule
Cell Combination & Voltage (V) & Observations \\
\midrule
Copper / Zinc & 0.94 & Lost a layer\\
&&Looked cleaner/shinier\\
Copper / Aluminum &1.52 & \\
Copper / Tin &0.39 & \\
\bottomrule
\end{tabular}
\end{table}

\vspace{0.5cm}
% pH
\noindent
\begin{table}[ht]
  \caption{pH of Vinegar}
\begin{tabular}{p{4cm} p{3.5cm} p{5cm}}
\toprule
Vinegar Sample & pH \\
\midrule
Vinegar & 2.4 \\
\bottomrule
\end{tabular}
\newpage
\end{table}

\section*{Analysis}

\subsection*{Step 1: Prediction of Cell Voltages}
To predict the cell potentials ($E_{\text{cell}}$) under non-standard conditions, the Nernst equation at $25^\circ\text{C}$ was utilized:
\[ E_{\text{cell}} = E^\circ_{\text{cell}} - \frac{0.0592\,\text{V}}{n}\log\left(Q\right) \]
where $Q$ is the reaction quotient, defined as:
\[ Q = \frac{[\text{M}^{n+}]}{[\text{H}^+]^2} \]

\subsubsection*{Calculation of Hydrogen Ion Concentration}
Given the measured pH of $2.4$ (Table 2), containing one significant digit in its fractional part, the hydronium concentration was determined as:
\[ [\text{H}^+] = 10^{-\text{pH}} = 10^{-2.4} = 4 \times 10^{-3}\text{ M} \]

\subsubsection*{Sample Calculation: Copper-Zinc Cell}
Using $E^\circ_{\text{cell}} = +0.76\text{ V}$, $n = 2$, and assuming an estimated trace metal ion concentration of $[\text{Zn}^{2+}] = 1.0 \times 10^{-5}\text{ M}$:
\begin{align*}
    E_{\text{cell}} &= E^\circ_{\text{cell}} - \frac{0.0592}{2}\log\left(\frac{[\text{Zn}^{2+}]}{[\text{H}^+]^2}\right) \\
    E_{\text{cell}} &= 0.76 - 0.0296\log\left(\frac{1.0 \times 10^{-5}}{(4 \times 10^{-3})^2}\right) \\
    E_{\text{cell}} &= 0.76 - 0.0296\log(0.625) \\
    E_{\text{cell}} &= 0.76 - 0.0296(-0.20) \\
    E_{\text{cell}} &= 0.76 + 0.0059 \\
    E_{\text{cell}} &\approx +0.77\text{ V}
\end{align*}

The predicted non-standard cell potential for the copper-zinc cell was determined to be +0.77 V.

\newpage

\subsection*{Step 2: Prediction of pH from Measured Voltage}
The Nernst equation was rearranged to solve for the pH of the electrolyte solution using the measured experimental cell potentials.

\subsubsection*{Mathematical Derivation}
Starting with the non-standard relation:
\[ E_{\text{cell}} = E^\circ_{\text{cell}} - \frac{0.0592}{2}\log[\text{M}^{2+}] + 0.0592\log[\text{H}^+] \]
Substituting $\text{pH} = -\log[\text{H}^+]$:
\[ E_{\text{cell}} = E^\circ_{\text{cell}} - 0.0296\log[\text{M}^{2+}] - 0.0592\text{ pH} \]
Rearranging to isolate pH yields:
\[ \text{pH} = \frac{E^\circ_{\text{cell}} - E_{\text{cell}} - 0.0296\log[\text{M}^{2+}]}{0.0592} \]

\subsubsection*{Sample Calculation: Copper-Zinc System}
Using the experimental value $E_{\text{cell}} = 0.94\text{ V}$ (Table 1) and assuming a trace metal concentration of $[\text{Zn}^{2+}] = 1.0 \times 10^{-10}\text{ M}$:
\begin{align*}
    \text{pH} &= \frac{0.76 - 0.94 - 0.0296\log(1.0 \times 10^{-10})}{0.0592} \\
    &= \frac{0.76 - 0.94 - 0.0296(-10)}{0.0592} \\
    &= \frac{0.76 - 0.94 + 0.296}{0.0592} \\
    &= \frac{0.116}{0.0592} \\
    &\approx \mathbf{1.96}
\end{align*}
The pH of the vinegar electrolyte calculated from the copper-zinc cell potential was determined to be 1.96.


\begin{table}[ht]
  \caption{Calculated vs. Measured Values}
\begin{adjustbox}{center}
\begin{tabular}{l c c c c}
\toprule
Cell System & Theoretical $E^\circ_{\text{cell}}$ (V) & Predicted $E_{\text{cell}}$ (V) & Measured $E_{\text{cell}}$ (V) & Calculated pH from $E_{\text{cell}}$ \\
\midrule
Copper / Zinc & +0.76 & +0.77 & 0.94 & 1.96 \\
Copper / Aluminum & +1.66 & +1.68 & 1.52 & 8.06 \\
Copper / Tin & +0.14 & +0.15 & 0.39 & -8.45 \\
\bottomrule
\end{tabular}
\end{adjustbox}
\end{table}

\section*{Discussion}

% Discussion questions
% o Question #1: Identify any significant sources of error in the lab. Identify at least two errors &
% explain their impact on the results observed.
% o Question #2: If there was an opportunity to re-do the lab, what changes would be made?
% o Question #3: Why did you choose to this lab for the DIY?

1. Two significant sources of systematic experimental error were identified during the trials. First, surface oxidation and impurities on the metal strips, noted in the qualitative observations of Table 1, acted as a physical barrier to electron transfer. This surface resistance restricted active electrochemical contact, thereby altering the cell potential from the theoretical Nernstian predictions. Second, the vinegar electrolyte was not refreshed between consecutive trials. This lack of replacement allowed for the progressive accumulation of metal ions and potential localized neutralization near the electrode surface, altering the pH of the boundary layers and shifting subsequent voltage measurements. Furthermore, because the electrode strips were rinsed only with water and not abrasively cleaned with steel wool or sandpaper between trials, pre-existing oxide layers remained on the surfaces, introducing additional resistance.

2. If the investigation were to be repeated, several procedural modifications would be implemented to optimize the reliability of the dataset. The metal electrodes would be thoroughly sanded with fine steel wool prior to each trial to completely remove passivating oxide layers. Additionally, fresh portions of the vinegar electrolyte would be utilized for each cell pairing, and direct pH measurements would be recorded for every trial to ensure consistency in the baseline hydronium ion concentration.

3. This specific laboratory experiment was chosen because of an interest in the practical construction and optimization of electrochemical cells. The exploration was motivated by a desire to investigate how changes in ion concentration under non-standard conditions shift cell potentials away from standard reduction potentials, and whether the resulting Nernstian deviations could be utilized to mathematically model and resolve unknown solution properties such as pH.

\section*{Conclusion}
The experimental cell potentials of non-standard vinegar-based cells were determined to be 0.94 V, 1.52 V, and 0.39 V for the copper-zinc, copper-aluminum, and copper-tin systems, respectively, and were successfully utilized to estimate electrolyte pH values. These values were systematically compared against theoretical cell potentials and direct pH measurements, demonstrating deviations that highlight the impact of surface oxidation and concentration shifts in weak acid systems.

\end{document}